% Guideline & Policy for Project Proposal
% Only one project proposal needs to be submitted per group. The idea behind this proposal is
% that I can give you feedback and suggestions to make you come up with a feasible plan for the
% project. You don’t have to do exactly what you say in your proposal and can even completely
% change your solving strategies afterwards if you want. But it’s important to have at least one
% reasonable plan/baseline to start from. The length should be 4 to 6 pages, not including
% appendices.


% Structure
% It should contain title, authors, abstract, introduction, related work, model/method, and
% experiments. In particular, all these sections should be complete (though they can be short)
% except for model/method and experiments.
\documentclass{article}

% Use final option for non-anonymous class submission
\usepackage[final]{neurips_2021}

% Recommended packages
\usepackage[utf8]{inputenc}
\usepackage[T1]{fontenc}
\usepackage{hyperref}
\usepackage{url}
\usepackage{booktabs}
\usepackage{amsfonts}
\usepackage{nicefrac}
\usepackage{microtype}
\usepackage{xcolor}
\usepackage{amsmath}
\usepackage{graphicx}

\title{Spring 2026 Analytics II Team Project Proposal Report}

% List the group and authors
\author{
  Group 1 \\[1ex]  % Adds Group 1 at the top
  Joshua Boyles\\
  Dave Julianto\\
  Hayden Burke\\
  Daniel Veiner\\
  Sebastian Hajdu
}

% Set the date manually (March 8, 2026)
%\date{March 8, 2026}

%\author{
%  First Author \\
%  Department of XYZ \\
%  University Name \\
%  City, Country \\
%  \texttt{author1@university.edu}
%  \And
%  Second Author \\
%  Department of XYZ \\
%  University Name \\
%  City, Country \\
%  \texttt{author2@university.edu}
%}


\begin{document}

\maketitle

% Add professor's name below title page (separate from authors)
\vspace{2em}  % Adds vertical space between the authors and professor's name
\begin{center}
  \textbf{Professor:} Dr. Kaixun Hua
\end{center}

% Abstract (15%)
% It should summarize the main idea of the project and its contributions. It should be
% understandable to anyone in the course.
\newpage
\begin{abstract}
Briefly summarize the problem, the proposed method, and the main findings.
Keep it concise (typically 150–250 words). Daniel volunteered to do this section.
\end{abstract}



% Introduction (15%)
% It should clearly state the problem being addressed and why it is important.
% From ChatGPT: Introduce the research problem, motivation, and contributions of the paper.
\section{Introduction}

This project aims to utilize historical ball bearing manufacturing data to predict 
defectiveness of manufactured parts and to understand which combinations of control 
settings typically lead to defective production. The goal is to create an interpretable 
model that accurately predicts whether a ball bearing will be classified as defective 
after several destructive and non-destructive tests, based on a set of predictor variables.

The dataset contains approximately 230,000 samples with 30 predictor variables. Each sample 
is labeled as defective (class 1) or non-defective (class 0), which is the target variable 
we aim to predict. By understanding the factors that contribute to defects, corrective 
action can be taken on controllable system inputs to reduce future defects. Reduced defect 
rates directly lower scrap, rework, resource, and labor costs in the manufacturing process.

However, this problem presents several challenges: the dataset is large and high-dimensional, 
the classes may be imbalanced (with relatively few defective samples), and manufacturing 
processes can involve complex, nonlinear interactions between variables. Additionally, 
interpretability is critical, as predictions must provide actionable insights for engineers 
and management.

In this project, we aim to:

\begin{enumerate}
    \item Develop an interpretable predictive model for classifying defective ball bearings.
    \item Identify and analyze key process variables that most strongly influence defect rates.
    \item Explore how the model’s insights can guide corrective actions to reduce defects and improve manufacturing efficiency.
\end{enumerate}

Together, these aims outline a plan to leverage historical manufacturing data for 
actionable defect prediction and process improvement.


% Related Work (15%)
% It should clearly distinguish what your proposed contribution from previous literature (e.g., a
% 1-2 sentence summary of at least 3 closely related papers).
\section{Related Work}
Discuss prior research and explain how your work differs from or extends it.



% Model/Method (25%)
% 1. It should contain the description of your research idea. You can use equations or sketches
% to help explain the idea.

% 2. It should contain a list of to-dos (having some time estimation could be helpful) that could
% in principle make you finish the project. Here’s an example:
% • Extended literature search and finishing related work section (3 hours)
% • Finding appropriate datasets and converting them to same format (2 hours)
% • Downloading xxx baseline and reproducing results of prior work (4 hours)
% • Implementing the proposed model (4 hours)
% • Running baselines on datasets and tweaking designs and or hyperparameters (4 hours)
% • Making figures and tables (2 hours)
% • Writing the project report (4 hours)

% You should almost always start from simple ideas and gradually add the complexity since
% a research idea typically has to FAIL many times until it works. It is highly likely that the
% project doesn’t go as planned. Therefore, with simple ideas, you could have the flexibility to
% tweak and a higher chance of obtaining some conclusions in the end.
\section{Model/Method}
Describe your model, algorithm, or methodology.

Example equation:
\begin{equation}
    y = f(x; \theta)
\end{equation}

The dataset exhibits a severe class imbalance which is one of the challenges being faced.
Out of the 229,788 observations, only 458 are defective, while 229,330 are non-defective.
As illustrated in figure 1 below, this disparity often causes classifiers to favor the majority class.
This can result in misleadingly high accuracy that masks manufacturing defects. Given that false negatives (missing a defect)
are significantly more costly than false positives, precision and recall are the primary metrics for evaluating this process. 

The modeling pipeline incorporates four class imbalance handling strategies:

\begin{enumerate}
    \item Random Oversampling.
    \item Random Undersampling
    \item Synthetic Minority Oversampling Technique (SMOTE)
    \item Class weighted Training
\end{enumerate}
%Experiments Design (20%)The model
% It should contain a list of to-do experiments (e.g., each experiment takes one subsection). One
% should clearly describe reasons of the experiments. One should also put placeholder figures and
% tables (with specifically designed axes, rows, columns, and so on). At last, one should write
% down the goal and expected outcome of specific experiments.

% Typically, in machine learning and related areas, one should have at least one main experiment
% on a convincingly large real-world dataset to show that your proposed model/method
% empirically works better than prior work. Moreover, one usually adds other experiments on
% real-world datasets to analyze how important a specific model/method design is, i.e., the so-
% called ablation study. One can also design special synthetic experiments to showcase the
% benefits of the proposed model/method under certain (possibly extreme) conditions.
\section{Experiments Design}
Describe the experimental setup, datasets, evaluation metrics, and results.
Include figures and tables as needed.

\bibliographystyle{plain}
\bibliography{references}

\end{document}
